% !TEX root = ../main.tex
\chapter{Unary Functions}
\label{ch:unary_functions}

The unary functions are the simplest op category: one input tensor,
one output tensor of the same shape, an element-wise scalar
transformation between them. \Lucid\ ships roughly 60 unary ops,
covering arithmetic, exponentials, logarithms, trigonometry,
hyperbolics, rounding, special functions, and the activations that
neural networks consume. This chapter enumerates them, gives each
its forward and backward, and documents the numerical
considerations that apply.

The category-level machinery was described in
\chref{ch:kernel_abstraction}: every unary op is a
\code{UnaryKernel<Op>} instantiation where \code{Op} declares the
scalar function and its derivative. The per-op files in
\code{lucid/\_C/ops/ufunc/} contain the \code{Op} structs and a
self-registering schema declaration. This chapter is the
mathematical exposition that those structs encode.

We organise the chapter by family: arithmetic
(\secref{sec:unary_functions:arith}), exponential
(\secref{sec:unary_functions:exp}), trigonometric
(\secref{sec:unary_functions:trig}), hyperbolic
(\secref{sec:unary_functions:hyper}), rounding
(\secref{sec:unary_functions:round}), special functions
(\secref{sec:unary_functions:special}), activations
(\secref{sec:unary_functions:activations}), and predicates
(\secref{sec:unary_functions:predicates}).

\section{Arithmetic Family}
\label{sec:unary_functions:arith}

The arithmetic family contains the unary operations that are
mathematical equivalents of arithmetic with a constant:

\begin{itemize}
  \item \code{neg(x) = -x}; backward: \code{grad\_in = -grad\_out}.
  \item \code{abs(x) = |x|}; backward: \code{grad\_in = sgn(x) *
    grad\_out}, undefined at 0 (we use 0 by convention).
  \item \code{sign(x) = sgn(x)}; backward: 0 everywhere.
  \item \code{square(x) = x\textasciicircum{}2}; backward: \code{grad\_in = 2*x*grad\_out}.
  \item \code{reciprocal(x) = 1/x}; backward: \code{grad\_in =
    -1/x\textasciicircum{}2 * grad\_out}.
  \item \code{sqrt(x)} returns $\sqrt{x}$; backward: $1/(2\sqrt{x})
    \cdot \text{grad\_out}$, undefined at 0.
  \item \code{rsqrt(x)} returns $1/\sqrt{x}$; backward:
    $-1/(2 x \sqrt{x}) \cdot \text{grad\_out}$.
  \item \code{cube(x) = x\textasciicircum{}3}; backward: \code{grad\_in = 3*x\textasciicircum{}2 *
    grad\_out}.
\end{itemize}

All but \code{sign} are differentiable; \code{sign} has zero
gradient by convention (since its mathematical derivative is a
Dirac delta).

\subsection{Convention at non-differentiable points}
\label{ssec:unary_functions:nondiff_conv}

\code{abs}, \code{sqrt}, and \code{rsqrt} have non-differentiable
points at zero. \Lucid's convention is the subgradient zero at the
non-differentiable point:

\begin{itemize}
  \item \code{abs'(0) := 0}
  \item \code{sqrt'(0) := 0} (rather than \(+\infty\))
  \item \code{rsqrt'(0)} raises a numerical warning; the input
    should never be exactly zero for an op whose forward already
    fails there.
\end{itemize}

The choice matches \refframework's convention and is what the
parity tests check.

\section{Exponential Family}
\label{sec:unary_functions:exp}

The exponential family:

\begin{itemize}
  \item \code{exp(x)} computes $e^x$; backward:
    \code{grad\_in = exp(x) * grad\_out} (saved tensor: output).
  \item \code{exp2(x)} computes $2^x$; backward:
    \code{grad\_in = ln(2) * exp2(x) * grad\_out}.
  \item \code{log(x)} computes $\ln(x)$; backward:
    \code{grad\_in = 1/x * grad\_out}.
  \item \code{log2(x)} computes $\log_2 x$; backward:
    \code{grad\_in = 1/(x ln(2)) * grad\_out}.
  \item \code{log10(x)} computes $\log_{10} x$; backward:
    analogous with $\ln(10)$.
  \item \code{log1p(x)} computes $\ln(1 + x)$; backward:
    \code{grad\_in = 1/(1+x) * grad\_out}. Numerically stable
    for small $x$.
  \item \code{expm1(x)} computes $e^x - 1$; backward:
    \code{grad\_in = (1 + expm1(x)) * grad\_out}. Numerically
    stable for small $x$.
\end{itemize}

\subsection{Numerical stability for small inputs}
\label{ssec:unary_functions:expm1_log1p}

\code{log1p} and \code{expm1} exist for the standard reason that
\code{log(1 + x)} and \code{exp(x) - 1} catastrophically cancel for
small \(x\) in floating point. The \code{log1p} primitive uses the
machine's hardware \code{log1p}-equivalent if available
(\code{vForce}'s \code{vvlog1pf} on CPU, an MLX primitive on GPU);
otherwise it falls back to a polynomial approximation that
maintains precision for \(|x| < 0.5\).

\section{Trigonometric Family}
\label{sec:unary_functions:trig}

Standard trigonometry, in radians:

\begin{itemize}
  \item \code{sin(x), cos(x), tan(x)}; backwards
    \code{cos(x) * grad\_out}, \code{-sin(x) * grad\_out}, \code{(1
    + tan(x)\textasciicircum{}2) * grad\_out}.
  \item \code{asin(x), acos(x), atan(x)}; backwards $1/\sqrt{1
    - x^2} \cdot \text{grad\_out}$, $-1/\sqrt{1 - x^2} \cdot
    \text{grad\_out}$, $1/(1 + x^2) \cdot \text{grad\_out}$.
\end{itemize}

The inverse trigonometric functions are undefined or
non-differentiable at \(|x| = 1\); we follow the IEEE 754
convention (returning \(\pm\pi/2\) for \code{asin} / \code{acos} at
the boundary and \(\pm\infty\)-equivalent for the derivative).

\section{Hyperbolic Family}
\label{sec:unary_functions:hyper}

Hyperbolic functions:

\begin{itemize}
  \item \code{sinh(x), cosh(x), tanh(x)}.
  \item \code{asinh(x), acosh(x), atanh(x)}.
\end{itemize}

\code{tanh}'s backward uses a saved-output optimisation: the
derivative is \(1 - \tanh^2(x)\), expressible in terms of the
output alone. The schema flags this and the autograd engine saves
the output rather than the input, halving the memory footprint
for \code{tanh}-heavy networks.

\section{Rounding Family}
\label{sec:unary_functions:round}

\begin{itemize}
  \item \code{floor(x)} computes $\lfloor x \rfloor$.
  \item \code{ceil(x)} computes $\lceil x \rceil$.
  \item \code{round(x)}: banker's rounding (round half to even),
    matching IEEE 754 default.
  \item \code{trunc(x)}: round toward zero.
  \item \code{frac(x)} computes $x - \mathrm{trunc}(x)$.
\end{itemize}

The rounding ops have zero gradient (they are piecewise constant
with measure-zero discontinuities at integers). The autograd
engine returns a zero gradient if backward is requested.

\subsection{IEEE 754 conformance}
\label{ssec:unary_functions:ieee754}

\Lucid's \code{round} matches IEEE 754: \code{round(0.5)} is 0
(round half to even), not 1. This differs from Python's
\code{round} when called on a float (Python also uses banker's
rounding) but matches \refframework. Parity tests check the
behaviour at half-integers.

\section{Special Functions}
\label{sec:unary_functions:special}

The special-function family is moved to its own subpackage
(\code{lucid.special}, \chref{ch:special_functions}) for
organisational reasons, but the underlying ops live in
\code{ops/ufunc/}. The list:

\begin{itemize}
  \item \code{erf(x)}: error function.
  \item \code{erfinv(x)}: inverse error function.
  \item \code{erfc(x)}: complementary error function.
  \item \code{lgamma(x), digamma(x), trigamma(x)}: gamma and its
    log-derivatives.
\end{itemize}

The forward implementations come from \Accelerate's \code{vForce}
on CPU and \MLX's special-function primitives on GPU. The
backwards are derived analytically and are encoded in the
\code{Op} structs.

\subsection{Numerical precision of erfinv}
\label{ssec:unary_functions:erfinv_precision}

\code{erfinv} is a difficult function to compute accurately near
\(|x| = 1\), where the function grows unboundedly. \vForce's
\code{vverfinvf} differs from scalar \code{erfinvf} by a few ULPs
in this region; we relax \code{gradcheck}'s tolerance for
\code{erfinv}-dependent ops accordingly. The behaviour is
documented per-op in the API reference.

\section{Activations}
\label{sec:unary_functions:activations}

The activation functions used by neural networks:

\begin{itemize}
  \item \code{relu(x)} computes $\max(0, x)$; backward:
    \code{grad\_in = (x > 0) * grad\_out}.
  \item \code{leaky\_relu(x, negative\_slope)}: \code{x} if
    $x > 0$, else \code{negative\_slope * x}.
  \item \code{prelu(x, alpha)}: \code{leaky\_relu} with learnable
    \code{alpha}.
  \item \code{elu(x)}: \code{x} if \code{x > 0}, else
    \code{alpha * (exp(x) - 1)}.
  \item \code{selu(x)}: scaled ELU (\code{alpha} and \code{scale}
    fixed to the self-normalising values).
  \item \code{gelu(x)}: Gaussian Error Linear Unit
    \cite{hendrycks_gelu_2016}. Exact
    formulation: $x \cdot 0.5 \cdot (1 + \text{erf}(x / \sqrt{2}))$.
    Approximate (tanh-based) formulation also available.
  \item \code{silu(x) = x * sigmoid(x)}: Sigmoid Linear Unit.
  \item \code{sigmoid(x) = 1 / (1 + exp(-x))}.
  \item \code{tanh(x)}: see hyperbolic family.
  \item \code{softplus(x)} computes $\ln(1 + e^x)$.
  \item \code{softsign(x) = x / (1 + |x|)}.
  \item \code{hardshrink(x, lambda)}: \code{x} if \code{|x| >
    lambda}, else 0.
  \item \code{hardtanh(x, min\_val, max\_val)}: clamp.
  \item \code{hardsigmoid(x)}: piecewise linear sigmoid.
  \item \code{hardswish(x) = x * hardsigmoid(x + 3)/6} (per the
    paper).
  \item \code{mish(x) = x * tanh(softplus(x))}.
\end{itemize}

\tabref{tab:unary_functions:activations} summarises the activation
family on one page. The mathematical forms are taken from the
respective papers; the \Lucid\ column gives the canonical
\code{lucid.nn.functional} entry point.

\begin{table}[t]
\centering
\footnotesize
\setlength{\tabcolsep}{4pt}
\begin{adjustbox}{max width=\textwidth, center}
\begin{tabular}{lllll}
\toprule
Name & Forward & Backward saved & Range & Lucid entry \\
\midrule
\code{relu}    & $\max(0, x)$              & input  & $[0, \infty)$ & \code{F.relu} \\
\code{leaky\_relu} & $x$ or $\alpha x$     & input  & $\mathbb{R}$  & \code{F.leaky\_relu} \\
\code{prelu}   & $x$ or $\alpha x$ (learn) & input, $\alpha$ & $\mathbb{R}$ & \code{F.prelu} \\
\code{elu}     & $x$ or $\alpha(e^x-1)$    & input, output & $(-\alpha, \infty)$ & \code{F.elu} \\
\code{selu}    & $\lambda \cdot \text{elu}(x)$ & input, output & $(-\lambda\alpha, \infty)$ & \code{F.selu} \\
\code{gelu}    & $x \cdot \Phi(x)$ (exact) & input  & $\mathbb{R}$  & \code{F.gelu} \\
\code{silu}    & $x \cdot \sigma(x)$       & input, output & $\mathbb{R}$ & \code{F.silu} \\
\code{sigmoid} & $1/(1 + e^{-x})$          & output & $(0, 1)$      & \code{F.sigmoid} \\
\code{tanh}    & hyperbolic tangent        & output & $(-1, 1)$     & \code{F.tanh} \\
\code{softplus} & $\ln(1 + e^x)$            & input  & $(0, \infty)$ & \code{F.softplus} \\
\code{softsign} & $x / (1 + |x|)$          & input  & $(-1, 1)$    & \code{F.softsign} \\
\code{hardtanh} & clamp$(x, l, u)$         & input  & $[l, u]$     & \code{F.hardtanh} \\
\code{hardsigmoid} & piecewise linear      & input  & $[0, 1]$     & \code{F.hardsigmoid} \\
\code{hardswish}   & $x \cdot \text{hsig}(x+3)/6$ & input & $\mathbb{R}$ & \code{F.hardswish} \\
\code{mish}    & $x \cdot \tanh(\text{softplus}(x))$ & input, output & $\mathbb{R}$ & \code{F.mish} \\
\bottomrule
\end{tabular}
\end{adjustbox}
\caption[Activation family.]{The activation functions exposed by
\Lucid. ``Backward saved'' lists which tensors the autograd engine
retains for the backward computation; functions whose derivative is
expressible in terms of the output (sigmoid, tanh, silu) save the
output rather than the input as a memory optimisation.}
\label{tab:unary_functions:activations}
\end{table}

\subsection{GELU exact vs approximate}
\label{ssec:unary_functions:gelu_modes}

\code{gelu} has two computational forms:

\begin{itemize}
  \item \textbf{Exact}: $x \cdot 0.5 \cdot (1 + \text{erf}(x /
    \sqrt{2}))$. Uses the error function; numerically precise.
  \item \textbf{Approximate}: $0.5\, x \cdot (1 + \tanh(\sqrt{2/\pi}
    \cdot (x + 0.044715\, x^3)))$. Polynomial approximation; slightly
    faster on hardware without fused erf.
\end{itemize}

The forms agree to within \(10^{-3}\) for \(|x| < 5\) and diverge
for larger \(|x|\). \Lucid\ defaults to the exact form. The user
can request the approximate form via
\code{nn.functional.gelu(x, approximate=`tanh')}; the parity tests
check both forms.

\subsection{The MPSGraph dispatch carve-out}
\label{ssec:unary_functions:mpsgraph_carveout}

Five activations dispatch through \MPSGraph\ rather than \MLX\ in
the 3.4 per-op carve-out
(\chref{ch:backend_dispatch}~Section~\ref{sec:backend_dispatch:mpsgraph_carveout}):
\code{gelu}, \code{silu}, \code{layer\_norm}-related activations,
\code{embedding\_backward}, \code{max\_pool\_backward}. The win is
substantial (30$\times$ for GELU, 28$\times$ for embedding
backward) and the numerical agreement with the \MLX\ path is
within tolerance.

\section{Predicates}
\label{sec:unary_functions:predicates}

The unary predicates return boolean tensors:

\begin{itemize}
  \item \code{isnan(x)}: true where \code{x} is NaN.
  \item \code{isinf(x)}: true where \code{x} is \(\pm\infty\).
  \item \code{isfinite(x)}: true where \code{x} is finite.
  \item \code{isclose(x, y, atol, rtol)}: elementwise tolerance
    comparison.
  \item \code{nan\_to\_num(x, nan, posinf, neginf)}: replace
    special values with finite numbers (this is technically not a
    predicate but is grouped with them).
\end{itemize}

Predicates are not differentiable; backward returns zero gradient
on the input.

\section{Numerical Analysis of Activations}
\label{sec:unary_functions:numerical_analysis}

The activation functions exhibit a range of numerical
characteristics that affect both forward stability and gradient
flow. This section catalogues the issues that have shaped \Lucid's
implementation choices.

\subsection{Vanishing gradient at saturation}
\label{ssec:unary_functions:vanishing_grad}

Sigmoid and tanh saturate at large $|x|$: $\sigma(x) \to 1$ as
$x \to \infty$ and $\sigma(x) \to 0$ as $x \to -\infty$, with the
derivative $\sigma(x)(1 - \sigma(x))$ dropping to floating-point
zero for $|x| \gtrsim 16$ in FP32 or $|x| \gtrsim 8$ in FP16. The
gradient through a saturated unit is effectively zero.

The practical implication for training: saturated activations stop
learning. The architectural responses include better
initialisation (Xavier, He), batch normalisation, and the move to
non-saturating activations (ReLU and its variants). \Lucid\
implements all three and the user can mix them freely.

\subsection{Dying ReLU}
\label{ssec:unary_functions:dying_relu}

ReLU has zero gradient for $x \leq 0$. A unit whose pre-activation
stays negative for all training inputs receives no gradient and
remains dead. Empirically, depending on initialisation and
learning rate, up to 40\,\% of units in a ReLU network may die
within a few hundred steps; the remaining units bear the
representation burden.

The architectural mitigations include leaky ReLU
($\alpha = 0.01$), PReLU (learnable $\alpha$), ELU, and the
smoother GELU/SiLU. The most common choice in modern Transformer
training is GELU; in classical vision networks ReLU still
dominates because the dying problem is masked by the high unit
count in convolutional layers.

\subsection{Saturating ranges and gradient flow}
\label{ssec:unary_functions:sat_ranges}

\tabref{tab:unary_functions:activation_grad} catalogues the
gradient behaviour at the extremes for each activation family.

\begin{table}[t]
\centering
\small
\begin{tabular}{llll}
\toprule
Activation & $x \to -\infty$ & $x \to +\infty$ & Smoothness \\
\midrule
relu      & $0$              & $1$              & $C^0$ \\
leaky\_relu & $\alpha$       & $1$              & $C^0$ \\
elu       & $\alpha e^x \to 0$ & $1$            & $C^1$ \\
gelu      & $0$              & $1$              & $C^\infty$ \\
silu      & $0$              & $1$              & $C^\infty$ \\
sigmoid   & $0$              & $0$              & $C^\infty$ \\
tanh      & $0$              & $0$              & $C^\infty$ \\
softplus  & $0$              & $1$              & $C^\infty$ \\
mish      & $0$              & $1$              & $C^\infty$ \\
\bottomrule
\end{tabular}
\caption[Gradient behaviour at extremes.]{Limiting gradient values
and smoothness class for each activation. The smooth ($C^\infty$)
activations (GELU, SiLU, sigmoid, tanh, softplus, mish) admit
higher-order derivatives, which matters for second-order
optimisers and for some interpretability techniques; the piecewise
($C^0$) activations (ReLU, leaky ReLU) do not.}
\label{tab:unary_functions:activation_grad}
\end{table}

\subsection{The exp / log family precision}
\label{ssec:unary_functions:exp_log_prec}

The exponential family deserves quantitative attention. In FP32,
\code{exp(x)} overflows at $x \approx 88.7$ and underflows at $x
\approx -103.3$. In FP16, overflow at $x \approx 11.1$,
underflow at $x \approx -17.3$. \code{log(x)} returns $-\infty$
at $x = 0$ and NaN for negative $x$.

\Lucid's implementations follow the hardware: FP32 \code{exp}
overflows to $+\infty$, FP16 \code{exp} on M-series GPU likewise.
Downstream operations that consume the inf (\code{div}, \code{mul}
by zero) propagate NaN; the autograd graph will then propagate the
NaN through backward, which the anomaly detector
(\code{lucid.autograd.set\_detect\_anomaly(True)}) catches and
reports with the originating op.

For workloads where the intermediate exp might overflow, the user
should compose through \code{logsumexp} or \code{softmax} rather
than the naive chain. The composite layer documents this in
\chref{ch:composite_ops}.

\section{Activation Selection in Practice}
\label{sec:unary_functions:selection}

A subjective but useful guide, based on the empirical training
behaviour we have observed in \Lucid's model zoo
(\partref{part:model_zoo}):

\begin{itemize}
  \item \textbf{Convolutional networks (ResNet, ConvNeXt, MobileNet).}
    Use ReLU by default. The dying-unit problem is well-tolerated
    in the high-channel-count regime, and the cheap gradient is
    worth the speed.
  \item \textbf{Transformer language models (BERT, GPT, T5).}
    Use GELU. The smoother gradient flows better through the
    deep stacks (12--96 layers) where ReLU's piecewise nature
    accumulates roughness.
  \item \textbf{Vision transformers (ViT, Swin).} GELU
    universally. The architectural inheritance from the Transformer
    family applies.
  \item \textbf{Object detection and segmentation (DETR, Mask
    R-CNN).} Mixed: ReLU in the convolutional backbones, GELU in
    the transformer heads.
  \item \textbf{Generative models (VAE, DDPM, NCSN).} ELU or SiLU.
    The smoother negative-input gradient helps the
    encoder/decoder paths.
  \item \textbf{Sequence models (LSTM, GRU).} Tanh and sigmoid as
    architecturally prescribed; the gates are the activations.
\end{itemize}

The recommendations are not laws; \Lucid\ supports all activations
on every model.

\section{Benchmarks}
\label{sec:unary_functions:benchmarks}

\tabref{tab:unary_functions:benchmark} reports per-op forward
throughput on M3~Max for several representative activations. The
inputs are FP32 tensors of $10^7$ elements (40\,MB), which fits in
the M3~Max's SLC and avoids DRAM-bandwidth confounds.

\begin{table}[t]
\centering
\small
\begin{tabular}{lrrl}
\toprule
Activation & Throughput (GFLOP/s) & Dispatch path & Notes \\
\midrule
\code{relu}    & 12\,400 & MLX        & near peak BW-bound \\
\code{sigmoid} & 9\,200  & MLX        & one transcendental per element \\
\code{tanh}    & 9\,500  & MLX        & similar to sigmoid \\
\code{gelu} (exact) & 11\,800 & \emph{MPSGraph} & 30$\times$ over MLX 3.3 baseline \\
\code{gelu} (approx) & 11\,200 & MLX        & polynomial approximation \\
\code{silu}    & 11\,400 & \emph{MPSGraph} & similar gain to GELU \\
\code{elu}     & 9\,000  & MLX        & one exp per element \\
\code{mish}    & 7\,800  & MLX        & two transcendentals per element \\
\bottomrule
\end{tabular}
\caption[Activation throughput.]{Per-op forward throughput on M3~Max
for representative activations, FP32, $10^7$ elements. The
``MPSGraph'' rows are the 3.4 per-op carve-out
(\chref{ch:mpsgraph_dispatch}); these activations dispatch through
\MPSGraph\ instead of \MLX\ because the \MPSGraph\ implementation
is substantially faster. The MLX-path activations are
bandwidth-bound and approach the SLC peak.}
\label{tab:unary_functions:benchmark}
\end{table}

\section{Summary and Forward References}
\label{sec:unary_functions:summary}

The unary functions are the most numerous op category, with around
60 distinct ops organised into eight families. Each is a
\code{UnaryKernel<Op>} instantiation; each declares a forward
function, a backward function (where differentiable), and an AMP
policy. Numerical stability is handled by op-specific tricks
(\code{log1p}, \code{expm1}, fused \code{logsumexp}-style
accumulation in FP32 for half-precision inputs).

The next chapter, \chref{ch:binary_functions}, treats the binary
ops --- elementwise operations with two inputs, where broadcasting
and type promotion enter. \chref{ch:reductions} treats reductions,
where the output shape differs from the input shape.

\begin{lucidnote}
For the user: every unary op is exposed both as a free function
(\code{lucid.relu(x)}) and as a Tensor method
(\code{x.relu()}). The two forms are aliases for the same
underlying implementation. Activations are also re-exported under
\code{lucid.nn.functional} and as stateful modules under
\code{lucid.nn} (\code{lucid.nn.ReLU()}).
\end{lucidnote}
